\section{Results --- EEG measurement-to-control}

\subsection{Setup (one paragraph)}
We freeze a trained backbone per LOSO fold and dump its penultimate latent $Z$ on the source subjects (the
held-out target subject is used for reporting only, never in selection, gating, or probing). On the frozen
$Z$ we (i) localize a candidate domain-rich, task-light subspace $V_D$ with the score-Fisher generalized
eigenproblem (\S2.2), (ii) delete it by the $M$-oblique direct-sum projector (\S2.3) and measure subject and
task decodability on the complement $RZ=(I-P_N)Z$ with linear and MLP probes, against a same-$k$ random-removal
control, and (iii) sweep a global LPC penalty (CE $+ \lambda \cdot I(Z;D|Y)$) to compare selective deletion with the
standard global objective. We study two representations: TSMNet (LogEig/SPD tangent latent, $d_z=210$) and
EEGNet (convolutional penultimate latent, $d_z=16$).

\subsection{TSMNet: subject leakage is high-dimensional and redundant; global LPC ``works'' only by collapse}
\textbf{In the high-dimensional TSMNet latent, the framework localizes subject-preferential low-rank directions,
but deleting them does not remove subject leakage.} Subject identity is almost perfectly decodable from
the full $Z$ (MLP 0.997, chance 0.125; Fig.~\ref{fig:tsmnet}A); the score-Fisher selector localizes a compact candidate
subspace (nDcand $\approx$ 3 of 210; Fig.~\ref{fig:tsmnet}D), yet deleting $V_D$ only \textbf{dents} subject decode (0.997$\to$0.96 MLP) ---
essentially matching same-$k$ random removal ($\approx$0.997; Fig.~\ref{fig:tsmnet}A,C) --- while preserving the task (0.75$\to$0.75;
Fig.~\ref{fig:tsmnet}B). Even deleting the full Fisher-discriminative subspace (the LDA cap of 7 directions for 8 source
subjects) leaves subject decode at 0.92--0.98: subject identity is \textbf{redundantly re-encoded} across the
high-dimensional latent, so \textbf{low-rank selective deletion is insufficient} \claimtag{C4}.

\textbf{Control is negative (global penalty).} The global LPC penalty \emph{appears} to remove leakage at high $\lambda$, but
the per-epoch curves (Fig.~\ref{fig:collapse}A--C) show this is a sharp, $\lambda$-tied \textbf{collapse of the representation to the
origin}: task cross-entropy rises to $\ln 4$ (chance), the feature norm and top singular value go to 0, and
the penalty itself goes to 0 (trivially satisfied by $Z\to0$). This is an objective-scaling pathology, not a
gradient explosion (the diagnostic encoder gradient is $\sim$10$\times$ smaller than in healthy low-$\lambda$ training and
never non-finite) and not a smooth geometric over-compression (the ``effective rank stays high'' is a
scale-invariant-metric artifact) \claimtag{C5}. The collapse is \emph{fixable} by a warm-up schedule (or, at $\lambda=1$, a
scale-normalized penalty), confirming its optimization origin --- but every collapse-free, task-preserving
variant leaves subject decode at the ERM level (Fig.~\ref{fig:collapse}D), i.e. \textbf{once the collapse is prevented, the global
penalty removes no leakage} \claimtag{C6}. In the tested TSMNet/2a setting, global LPC's apparent de-domaining is a
collapse artifact.

\begin{figure}[t]
\centering
\includegraphics[width=\linewidth]{fig3_lpc_collapse_mechanism.pdf}
\caption{\textbf{Global LPC removes subject leakage only by collapsing the representation (TSMNet/2a).}
Raw global LPC removes subject decodability only by collapsing the latent to the origin; when this collapse
path is blocked or avoided, task performance recovers but subject leakage returns. Per-epoch medians:
(A) task cross-entropy rises to $\ln 4$ (chance) for $\lambda\ge1$; (B) feature norm $\|Z\|\to0$ (collapse to
origin); (C) $\lambda\cdot$LPC penalty $\to0$ (trivially satisfied by $Z\to0$). This is an objective-scaling /
training-basin pathology, \emph{not} a gradient explosion (the diagnostic encoder gradient is $\sim$10$\times$
smaller than in healthy low-$\lambda$ training and never non-finite; the stable effective rank is a
scale-invariant-metric artifact). (D) Raw LPC at $\lambda{=}1$ drives source task and subject decode to chance
($Z\to0$); warm-up or scale-normalization prevents the collapse, restoring task accuracy, but subject
leakage returns to the ERM level.}
\label{fig:collapse}
\end{figure}

\begin{figure}[t]
\centering
\includegraphics[width=\linewidth]{fig4_tsmnet_redundant_leakage.pdf}
\caption{\textbf{TSMNet subject leakage is high-dimensional and redundant} (ERM; $d_z{=}210$; $n{=}27$
folds$\times$seeds). (A) Subject decode for the full $Z$, the complement $RZ$ after deleting $V_D$, the
deleted $V_D$ itself, and a same-$k$ random removal (linear and MLP; chance dashed): deleting $V_D$ only
\emph{dents} subject decode and matches random removal. (B) Task decode is preserved. (C) Subject decode
\emph{removed} by informed $V_D$ vs same-$k$ random is small and comparable --- the deletion is not
selective. (D) Candidate rank nDcand ($\approx$3 of 210) is compact, but low-rank deletion is nonetheless
insufficient because subject identity is redundantly re-encoded across the high-dimensional latent.}
\label{fig:tsmnet}
\end{figure}

\subsection{EEGNet: low-rank deletion linearly reduces leakage, but removal yields no generalization benefit}
\textbf{EEGNet shows that the negative TSMNet result is not because score-Fisher deletion can never remove
leakage: in a compact latent the same diagnostic deletion removes much more subject information --- yet this
removal does not translate into target-domain improvement.} Deleting $V_D$ (nDcand $\approx$ 5 of 16) preserves the
task (0.64$\to$0.64) and \textbf{selectively removes} subject decode far below same-$k$ random removal: linear
0.82$\to$0.35 (random 0.73), MLP 0.88$\to$0.54 (random 0.81; Fig.~\ref{fig:eegnet}A--C). The informed-vs-random selectivity is
$\sim$0.35--0.55, roughly an order of magnitude larger than on TSMNet (0.04--0.08), and same-$k$ random deletion
does \emph{not} reproduce it --- genuine subspace selectivity, not an artifact of deleting a larger fraction of a
small latent \claimtag{C7}. The removal is \textbf{partial}: a substantial nonlinear residual survives (RZ MLP 0.54, well
above chance 0.125), so we report it as \emph{linearly reducible with a persistent nonlinear residual}, not as
elimination \claimtag{C8}.

\textbf{But removable $\ne$ beneficial.} Sweeping the global LPC penalty on EEGNet reduces subject decode
monotonically (0.89$\to$0.17) \emph{without} collapse (the feature norm never goes to 0; Fig.~\ref{fig:eegnet}D), yet \textbf{mean} LOSO
target accuracy is flat-to-worse (0.43$\to$0.39 across $\lambda$; paired-t worse at every $\lambda$, p$\le$0.001 at $\lambda\ge$1) and
uncorrelated with leakage reduction (Pearson $-$0.14, n.s.). Removing the measured leakage --- by selective
deletion or by a collapse-free global penalty --- does \textbf{not} improve cross-subject generalization on
EEGNet/2a \claimtag{C9}. \emph{(We report the mean LOSO accuracy, the standard DG metric; the median is $\sim$flat and would
obscure the paired decline.)}

\begin{figure}[t]
\centering
\includegraphics[width=\linewidth]{fig5_eegnet_contrast.pdf}
\caption{\textbf{EEGNet contrast: low-rank deletion linearly reduces leakage, but removal yields no DG gain}
(ERM $d_z{=}16$ and its global-LPC sweep). (A--C) Same axes as Fig.~\ref{fig:tsmnet}: deleting $V_D$
\emph{selectively} removes subject decode far below same-$k$ random removal (linear 0.82$\to$0.35, MLP
0.88$\to$0.54) while preserving the task --- an order-of-magnitude larger informed-vs-random selectivity
than TSMNet; a nonlinear residual remains (RZ MLP 0.54 $\gg$ chance). (D) Global LPC reduces subject decode
(0.89$\to$0.17) \emph{without} collapse (feature norm never $\to0$), yet \emph{mean} LOSO target accuracy is
flat-to-worse ($0.43\to0.39$, paired-$t$ worse, $p\le0.001$ at $\lambda\ge1$) --- removing the leakage does not improve target
generalization. Curves show the mean over folds$\times$seeds ($n{=}27$).}
\label{fig:eegnet}
\end{figure}

% Auto-generated by tos_cmi/paper/scripts/build_table1_summary.py.
% Do not edit values manually; edit the source builder or artifact loaders.
% Requires: \usepackage{booktabs,tabularx,makecell}
\begin{table}[t]\centering\footnotesize
\caption{Measurement-to-control summary across EEG representations (BCI-IV-2a, LOSO). Target DG = MEAN LOSO target balanced accuracy over folds x seeds (paired-t vs lambda=0). Deletion decisions are diagnostic unless the exact certificate scope is reached.}
\label{tab:summary}
\begin{tabularx}{\linewidth}{l X X}
\toprule
Property & \textbf{TSMNet} & \textbf{EEGNet} \\
\midrule
Latent dim & \makecell[tl]{210} & \makecell[tl]{16} \\
Subject decode (ERM) & \makecell[tl]{0.997} & \makecell[tl]{0.881} \\
Low-rank deletion (subject; task) & \makecell[tl]{dents only: subj 1.00 $\to$ 0.96 (= random 1.00); task kept} & \makecell[tl]{removes: subj linear 0.82 $\to$ 0.35 / MLP 0.88 $\to$ 0.54 ($\gg$ random 0.81); task kept} \\
Raw LPC outcome & \makecell[tl]{leakage falls via collapse (feature-norm $\to$ 0), task $\to$ chance} & \makecell[tl]{no collapse: leakage falls 0.89 $\to$ 0.17, source task degrades gradually} \\
Collapse-free LPC outcome & \makecell[tl]{task restored, leakage returns $\to$ ERM (subj ~0.997)} & \makecell[tl]{raw LPC already collapse-free on this latent} \\
DG gain under task-preserving control & \makecell[tl]{not achievable (no task-preserving leakage reduction: raw LPC collapses; collapse-free LPC removes none)} & \makecell[tl]{none: mean target 0.43 $\to$ 0.39, paired-t p<0.001 as leakage falls 0.89 $\to$ 0.17} \\
Certified / diagnostic decision & \makecell[tl]{abstain / low-rank deletion insufficient} & \makecell[tl]{diagnostic deletion removes leakage but no DG gain} \\
\bottomrule
\end{tabularx}
\end{table}


\subsection{Stronger linear erasure does not close the control gap}
\textbf{The measurement-to-control gap is not an artifact of a weak eraser: a stronger, closed-form optimal
linear eraser removes more subject information than the TOS deletion, yet the gap remains} (Table~\ref{tab:erasure}).
On both backbones LEACE --- the closed-form optimal linear eraser --- drives \emph{linearly} decodable
subject identity to chance ($0.12$) and reduces subject decodability more than the TOS $V_D$ deletion at
comparable task cost (MLP residual TSMNet $0.96\to0.74$, EEGNet $0.55\to0.39$); the rank-limited erasers
(RLACE, TOS) by contrast fail on the redundant high-dimensional TSMNet latent ($0.86$--$0.96$) and only work
on the compact EEGNet latent. So the score-Fisher \emph{projection is not} the methodological advantage~\claimtag{C7}.
The stronger eraser instead sharpens the phenomenon: after \emph{optimal} linear erasure a large
\textbf{nonlinearly recoverable subject residual} remains under held-out MLP probes on TSMNet ($0.74\gg$
chance $0.125$) versus a smaller one on EEGNet ($0.39$) --- consistent with nonlinear/redundant subject
coding that linear erasure cannot reach. INLP drives \emph{linear} subject decodability near chance (a nonlinear residual persists on TSMNet) only by sacrificing
task information (task $0.75\to0.55$ TSMNet, $0.64\to0.25$ EEGNet): unconstrained erasure is not task-safe.
Crucially, these controls establish that the source-side removal result is not limited by a weak eraser; whether
source-side removal improves \emph{target} accuracy is a separate question, which the global-LPC sweep (\S4.3)
answers in the negative for this setting (a frozen-erasure target-deployment test is left to future work). Successful
erasure is therefore not the right endpoint; the relevant question is whether the intervention is task-safe
\emph{and} beneficial for source-to-target generalization~\claimtag{C9,C10}.

% Auto-generated by tos_cmi/paper/scripts/build_table2_erasure.py. Do not hand-edit values.
% Requires: \usepackage{booktabs,tabularx}
\begin{table}[t]\centering\footnotesize
\caption{Linear concept-erasure controls show the projector is not the contribution (BCI-IV-2a, frozen latents; subject decode after erasing $D$, lower is more removed; chance subj.\ $0.125$, task $0.250$; task linear before erasure $0.75$/$0.64$). LEACE (optimal closed-form linear erasure) removes more subject information than the TOS deletion at equal task cost, yet a large nonlinearly-recoverable MLP residual remains (TSMNet $0.74$ vs EEGNet $0.39$); INLP erases only by destroying the task; rank-limited erasers (RLACE, TOS) fail on the redundant high-dimensional TSMNet latent.}
\label{tab:erasure}
\begin{tabularx}{\linewidth}{l l c c c X}
\toprule
Backbone & Method & Subj.\ lin. & Subj.\ MLP & Task lin. & Interpretation \\
\midrule
TSMNet ($d_z{=}210$) & (none) & 1.00 & 1.00 & 0.75 & reference (no erasure) \\
 & LEACE & 0.12 & 0.74 & 0.74 & optimal linear eraser; strongest linear removal at near-zero task cost \\
 & RLACE & 0.86 & 0.94 & 0.74 & rank-$k$ adversarial; rank-limited erasure \\
 & INLP & 0.14 & 0.46 & 0.55 & erases but destroys the task (over-erasure) \\
 & TOS ($V_D$) & 0.92 & 0.96 & 0.75 & low-rank score-Fisher deletion (rank-limited) \\
 & random-$k$ & 1.00 & 1.00 & 0.74 & same-$k$ random removal (control) \\
\midrule
EEGNet ($d_z{=}16$) & (none) & 0.82 & 0.88 & 0.64 & reference (no erasure) \\
 & LEACE & 0.12 & 0.39 & 0.61 & optimal linear eraser; strongest linear removal at near-zero task cost \\
 & RLACE & 0.29 & 0.49 & 0.61 & rank-$k$ adversarial; rank-limited erasure \\
 & INLP & 0.12 & 0.12 & 0.25 & erases but destroys the task (over-erasure) \\
 & TOS ($V_D$) & 0.36 & 0.54 & 0.64 & low-rank score-Fisher deletion (rank-limited) \\
 & random-$k$ & 0.73 & 0.81 & 0.62 & same-$k$ random removal (control) \\
\bottomrule
\end{tabularx}
\end{table}


\subsection{Removability is capacity-mediated, not architecture-specific}
\textbf{A seed-0 architecture$\times$latent-dimension factorial indicates that the TSMNet/EEGNet contrast is
driven largely by latent \emph{dimension} rather than the SPD-vs-convolutional representation type} (multi-seed
hardening, and a dimension-matched convolutional cell at $d_z{=}210$, are in progress). We vary the latent dimension
of each backbone (TSMNet tangent dimension via the SPD size; EEGNet penultimate width) and measure the
nonlinear residual after optimal linear erasure (Fig.~\ref{fig:capacity}). The residual rises monotonically
with latent dimension \emph{within both architectures}, and at \emph{matched} dimension the two are nearly
indistinguishable (e.g.\ $d_z\!\approx\!16$--$21$: $0.40$ vs $0.40$; $d_z\!\approx\!32$--$36$: $0.50$ vs
$0.50$; $d_z\!\approx\!105$--$128$: $0.65$ vs $0.61$). The earlier TSMNet-$0.74$-vs-EEGNet-$0.39$ gap is
thus largely attributable to capacity ($210$ vs $16$) rather than representation type: larger latents support
redundant, nonlinearly-recoverable subject coding that survives optimal linear erasure. This addresses the
dimension-vs-type confound that the two-point comparison alone could not, pending multi-seed confirmation (\S6.3).

\begin{figure}[t]
\centering
\includegraphics[width=\linewidth]{fig6_capacity_factorial.pdf}
\caption{\textbf{Removability is capacity-mediated.} Removability signals vs latent dimension $d_z$ for
TSMNet (SPD/LogEig) and EEGNet (conv). (A) full-$Z$ subject decode rises with $d_z$; (B) the MLP residual
after \emph{optimal} linear erasure (LEACE) rises with $d_z$ and the two architectures nearly coincide at
matched $d_z$ --- the contrast is capacity-mediated, not architecture-type; (C) task decode is preserved
after LEACE across $d_z$; (D) rank-limited deletion (TOS) degrades toward random-$k$ as $d_z$ grows. (Seed 0;
multi-seed runs in progress.)}
\label{fig:capacity}
\end{figure}

\subsection{Unified interpretation}
Across both representations, conditional domain leakage is a \textbf{measurable} property of the EEG latent that
the score-Fisher diagnostics localize, but it is \textbf{not a sufficient control target} for cross-subject
generalization (Table~\ref{tab:summary}). Removing it can (i) collapse the representation (TSMNet global LPC), (ii) leave
redundant high-dimensional leakage essentially intact (TSMNet low-rank deletion), or (iii) genuinely
remove leakage yet not improve transfer (EEGNet). Whether the leakage even \emph{sits} in a low-rank removable
subspace appears \textbf{capacity-mediated} (a seed-0 factorial, \S4.5): a higher-capacity latent redundantly
re-encodes subject identity outside any low-rank subspace, while a compact latent cannot. The correct operating principle is therefore a \textbf{certified intervention with refusal}: the
framework should delete only when a source-risk certificate permits, and abstain otherwise --- which on
TSMNet means abstain (deletion insufficient) and on EEGNet means a diagnostic deletion that removes leakage
but is not claimed to improve generalization \claimtag{C10}.

\textbf{In sum,} on these representations conditional domain leakage is measurable and sometimes (partially) removable, yet neither its measurement nor its removal demonstrates a cross-subject generalization benefit.

\emph{Limitation flagged here (full text in \S6.3): the two backbones differ in both architecture and latent
dimension; we disentangle these with a capacity factorial (\S4.5; seed-0, multi-seed in progress).}
